\section*{Chapitre \ref{ch:b1:quantum-introduction} --- Introduction à la physique quantique}

\begin{solution}{exo:b1:quantum-introduction:1}
$h\nu$~: \qty{100}{MHz}~: $6.6 \times 10^{-26}$~J $= \qty{4.1e-7}{eV}$~; \qty{500}{nm}~:
$1240/500 = \qty{2.5}{eV}$~; \qty{0.10}{nm}~: \qty{12.4}{keV}~; \qty{1}{MeV}~: $\lambda =
1240/10^6 = \qty{1.2e-3}{nm} = \qty{1.2}{pm}$. Laser~: $1240/633 = \qty{1.96}{eV} =
\qty{3.1e-19}{J}$, donc $10^{-3}/3.1 \times 10^{-19} = 3.2 \times 10^{15}$ photons par
seconde.
\end{solution}

\begin{solution}{exo:b1:quantum-introduction:2}
$\lambda_0 = 1240/2.3 = \qty{540}{nm}$. À \qty{400}{nm}, $h\nu = \qty{3.1}{eV}$~:
$E_{k,\max} = \qty{0.8}{eV}$, $V_s = \qty{0.8}{V}$. Doubler l'intensité
double le courant et laisse $V_s$ inchangé. À \qty{600}{nm}
(\qty{2.07}{eV} $< W$)~: rien, quelle que soit l'intensité.
\end{solution}

\begin{solution}{exo:b1:quantum-introduction:3}
Électron~: $1.23/\sqrt{100} = \qty{0.12}{nm}$~; à \qty{1}{keV}~: \qty{0.039}{nm}.
Neutron~: $E = \qty{4.0e-21}{J}$, $p = \sqrt{2m_nE} = \qty{3.7e-24}{kg.m/s}$,
$\lambda = h/p = \qty{0.18}{nm}$. Bille~: $6.6 \times 10^{-34}/10^{-3} = \qty{7e-31}{m}$.
Les trois premières ont un $\lambda$ de l'ordre de la distance interatomique et
sont diffractées (cristallographie électronique, neutronique, par rayons X)~;
la bille, jamais.
\end{solution}

\begin{solution}{exo:b1:quantum-introduction:4}
$E_1 = h^2/8m_eL^2 = 4.39 \times 10^{-67}/(8 \times 9.11 \times 10^{-31} \times 10^{-18}) =
\qty{6.0e-20}{J} = \qty{0.38}{eV}$, $E_2 = 4E_1 = \qty{1.5}{eV}$~; photon de
\qty{1.13}{eV}, $\lambda = 1240/1.13 = \qty{1.1}{\micro m}$ (proche infrarouge). Proton
dans \qty{5}{fm}~: $4.39 \times 10^{-67}/(8 \times 1.67 \times 10^{-27} \times 25 \times 10^{-30}) =
\qty{1.3e-12}{J} = \qty{8}{MeV}$ --- l'échelle des énergies nucléaires, par le
seul confinement.
\end{solution}

\begin{solution}{exo:b1:quantum-introduction:5}
Pente $(2.05 - 0.40)/(4.0 \times 10^{14}) = \qty{4.1e-15}{V.s}$, $h = e \times 4.1 \times
10^{-15} = \qty{6.6e-34}{J.s}$. $W = h\nu - eV_s = 2.48 - 0.40 = \qty{2.1}{eV}$,
seuil $\lambda_0 = 1240/2.1 = \qty{600}{nm}$~: le césium.
\end{solution}

\begin{solution}{exo:b1:quantum-introduction:6}
$\Delta\lambda = \qty{2.43}{pm}$~: $\lambda' = \qty{12.4}{pm}$. Avant~: $1240/0.010 =
\qty{124}{keV}$~; après~: $1240/0.0124 = \qty{100}{keV}$~; l'électron emporte
\qty{24}{keV}. Pour la lumière, $\Delta\lambda/\lambda = 2.4 \times 10^{-12}/5 \times 10^{-7} = 5 \times
10^{-6}$~: immesurable en 1923 --- les rayons X portent le décalage à $25\%$.
\end{solution}

\begin{solution}{exo:b1:quantum-introduction:7}
$p = \sqrt{2 \times 9.11 \times 10^{-31} \times 5 \times 10^4 \times 1.6 \times 10^{-19}} =
\qty{1.2e-22}{kg.m/s}$, $\lambda = \qty{5.5}{pm}$ (la relativité la corrige en
\qty{5.4}{pm})~; interfrange $\lambda D/d = 5.5 \times 10^{-12} \times 0.35/2 \times 10^{-6}
= \qty{1}{\micro m}$~: invisible à l'\omterm{def:b1:optical-instruments:eye}{œil}, d'où la \omterm{def:b1:mirrors-thin-lenses:lens}{lentille} grossissante. Dix
électrons~: dix points, apparemment au hasard~; $10^5$~: les franges. Chaque
électron arrive en un point~; c'est la \emph{probabilité} d'y arriver qui suit
l'onde des deux fentes --- l'onde d'un seul électron passe par les deux.
\end{solution}

\begin{solution}{exo:b1:quantum-introduction:8}
Atome~: $\Delta p \geq \hbar/2\Delta x = \qty{5.3e-25}{kg.m/s}$, $E \sim \Delta p^2/2m_e =
\qty{1.5e-19}{J} \approx \qty{1}{eV}$~: le bon ordre. Noyau~: $\Delta p \geq
\qty{1.1e-20}{kg.m/s}$~; $\Delta p^2/2m_e$ vaudrait \qty{400}{MeV}, bien au-delà de
$m_ec^2$, donc $E \approx \Delta p\,c = \qty{3e-12}{J} = \qty{20}{MeV}$~: aucune force
nucléaire ne retiendrait un électron aussi énergétique --- les noyaux ne
contiennent pas d'électrons (ceux du rayonnement bêta sont créés à l'émission).
Grain~: $\Delta v \geq
\hbar/2m\Delta x = 1.05 \times 10^{-34}/(2 \times 10^{-9} \times 10^{-6}) = \qty{5e-20}{m/s}$~:
sans conséquence.
\end{solution}

\begin{solution}{exo:b1:quantum-introduction:9}
(a) $m_ev^2/r = e^2/4\pi\varepsilon_0r^2$~: $v = \sqrt{e^2/4\pi\varepsilon_0m_er}$, $E =
\tfrac12m_ev^2 - e^2/4\pi\varepsilon_0r = -e^2/8\pi\varepsilon_0r$. (b) $m_evr = n\hbar$ et
$m_ev^2r = e^2/4\pi\varepsilon_0$~: diviser le carré de la première par la seconde~:
$m_er = n^2\hbar^2 \cdot 4\pi\varepsilon_0/e^2$, $r_n = n^2a_0$, et $E_n = -e^2/8\pi\varepsilon_0n^2a_0
= E_1/n^2$ avec $a_0 = \qty{52.9}{pm}$, $E_1 = \qty{-13.6}{eV}$. (c) $v_1 =
\hbar/m_ea_0 = \qty{2.19e6}{m/s}$, $v_1/c = 1/137$. (d) $2 \to 1$~: $\tfrac34 \times
13.6 = \qty{10.2}{eV}$, \qty{122}{nm}~; $3 \to 2$~: $(\tfrac14 - \tfrac19) \times 13.6 =
\qty{1.89}{eV}$, \qty{656}{nm}. (e) $\lambda_n = h/m_ev_n = nh/m_ev_1 = 2\pi na_0$
(puisque $m_ev_1a_0 = \hbar$), et $2\pi r_n = 2\pi n^2a_0 = n\lambda_n$.
\end{solution}

\begin{solution}{exo:b1:quantum-introduction:10}
(a) Remplacer $e^2$ par $Ze^2$~: $a_0 \to a_0/Z$, $E_1 \to Z^2E_1$. He$^+$~:
\qty{54.4}{eV}. U$^{91+}$~: $92^2 \times 13.6 = \qty{115}{keV}$ --- mais $v_1 = Z\alpha c
= 0.67c$~: le modèle non relativiste n'y est qu'indicatif. (b) $a_0/207
= \qty{256}{fm}$, $E_1 = 207 \times 13.6 = \qty{2.8}{keV}$~; Lyman $\alpha$~:
\qty{2.1}{keV}, $\lambda = \qty{0.59}{nm}$ (rayons X). L'orbite du muon ne fait que
$300$ rayons du proton~: ses niveaux se déplacent sensiblement avec la taille
du proton, et c'est ainsi qu'on l'a remesurée (2010). (c) Les deux
particules tournent autour de leur centre commun~: la \omterm{prop:b1:systems-of-points:twobody}{masse réduite} $m_e/2$
divise par deux toutes les énergies~: $E_1 = \qty{-6.8}{eV}$, Lyman $\alpha$ à \qty{5.1}{eV}, \qty{243}{nm}.
\end{solution}

\begin{solution}{exo:b1:quantum-introduction:11}
(a) $h^2/8m_e = \qty{6.02e-38}{J.m^2}$. $R = \qty{2}{nm}$~: $E_e = 6.02 \times 10^{-38}/
(0.13 \times 4 \times 10^{-18}) = \qty{0.72}{eV}$, $E_h = \qty{0.21}{eV}$, au total $1.74 +
0.93 = \qty{2.67}{eV}$, \qty{464}{nm}~: bleu. \qty{3}{nm}~: confinement $\times4/9$,
\qty{2.15}{eV}, \qty{577}{nm}~: jaune. \qty{4}{nm}~: $\times1/4$, \qty{1.97}{eV},
\qty{629}{nm}~: rouge. (b) \qty{520}{nm} font \qty{2.38}{eV}~: confinement \qty{0.65}{eV}
$= 0.93 \times (2/R)^2$, $R = \qty{2.4}{nm}$. (c) Sans confinement~: $E_g$ seul,
\qty{713}{nm}. (d) $\delta E = -2 \times 0.1 \times 0.41 = \qty{-0.08}{eV}$, soit $4\%$ de
\qty{2.15}{eV}~: environ \qty{20}{nm} de dispersion --- des couleurs pures
exigent des lots monodisperses.
\end{solution}

\begin{solution}{exo:b1:quantum-introduction:12}
(a) $E(\Delta x) = \hbar^2/8m\Delta x^2 + \tfrac12m\omega^2\Delta x^2$~; dérivée nulle en
$\Delta x^2 = \hbar/2m\omega$, où les deux termes valent $\hbar\omega/4$~: $E_{\min} =
\hbar\omega/2$. (b) $\tfrac12 \times 1.055 \times 10^{-34} \times 8.3 \times 10^{14} = \qty{4.4e-20}{J}
= \qty{0.27}{eV}$~; $T \sim \hbar\omega/k_B = \qty{6300}{K}$ --- à température ambiante
la vibration est gelée dans son \omterm{prop:b1:quantum-introduction:box}{état fondamental}. (c) $\Delta p = \hbar/2\Delta x =
\qty{1.05e-24}{kg.m/s}$, $E = \Delta p^2/2m_{\mathrm{He}} = 1.1 \times 10^{-48}/1.33 \times
10^{-26} = \qty{8e-23}{J} = \qty{0.5}{meV}$~: la moitié de l'énergie de liaison ---
les atomes ne tiennent pas assez tranquilles pour cristalliser~; l'hélium
reste liquide jusqu'au zéro absolu tant qu'on ne le comprime pas
(\qty{25}{bar}). (d) $E_1 = \pi^2\hbar^2/2mL^2$ contre
$\hbar^2/2mL^2$~: dix fois plus grand, et il le faut --- la borne est une
minoration, et l'état de la boîte a $\Delta x \approx 0.18L$, non $L/2$.
\end{solution}

\begin{solution}{pb:b1:quantum-introduction:1}
\textbf{1.} La lumière arrache à la cathode des électrons d'énergies
diverses~; portée à $-V_s$, l'anode les repousse et seuls arrivent ceux dont
$E_k > eV_s$. Le courant s'annule lorsque $eV_s = E_{k,\max}$.

\textbf{2.} $eV_s = h\nu - W$~: pente $h/e$, ordonnée à l'origine $-W/e$~; $V_s = 0$ au
seuil $\nu_0 = W/h$.

\textbf{3.} Points extrêmes~: $(2.36 - 0.30)/(5.0 \times 10^{14}) = \qty{4.12e-15}{V.s}$,
$h = 1.602 \times 10^{-19} \times 4.12 \times 10^{-15} = \qty{6.60e-34}{J.s}$ (une régression
sur les cinq points donne $6.62$)~: à $0.5\%$ de $6.626 \times 10^{-34}$.

\textbf{4.} $W = h\nu - eV_s = 2.47 - 0.30 = \qty{2.2}{eV}$~; $\nu_0 = W/h =
\qty{5.3e14}{Hz}$, $\lambda_0 = 1240/2.2 = \qty{560}{nm}$. Un pointeur rouge
(\qty{650}{nm}, \qty{1.9}{eV}) n'arrache rien.

\textbf{5.} $V_s$ inchangé, courant doublé~: deux fois plus de photons, chacun
de même énergie. Une onde donnerait à chaque électron plus d'énergie à mesure
que l'intensité croît, et un seuil en intensité plutôt qu'en fréquence.

\textbf{6.} $10^{-6} \times 10^{-20} = \qty{1e-26}{W}$ sur l'atome~; \qty{2.2}{eV} $=
\qty{3.5e-19}{J}$ demandent $3.5 \times 10^7$~s --- un an. Or l'émission s'observe
en quelques nanosecondes~: l'énergie arrive par quanta entiers.

\textbf{7.} $m_ev^2/r = e^2/4\pi\varepsilon_0r^2$~: $v = \sqrt{e^2/4\pi\varepsilon_0m_er}$~; $E =
\tfrac12m_ev^2 - e^2/4\pi\varepsilon_0r = -e^2/8\pi\varepsilon_0r$.

\textbf{8.} $(m_evr)^2 = n^2\hbar^2$ et $m_ev^2r = e^2/4\pi\varepsilon_0$~: $r_n = n^2 \cdot
4\pi\varepsilon_0\hbar^2/m_ee^2 = n^2a_0$, $a_0 = (1.055 \times 10^{-34})^2/(8.99 \times 10^9
\times 9.11 \times 10^{-31} \times 2.57 \times 10^{-38}) = \qty{52.9}{pm}$.

\textbf{9.} $E_n = -e^2/8\pi\varepsilon_0r_n = -m_ee^4/2(4\pi\varepsilon_0)^2\hbar^2n^2$~; $E_1 =
-2.31 \times 10^{-28}/(2 \times 5.29 \times 10^{-11}) = \qty{-2.18e-18}{J} = \qty{-13.6}{eV}$.

\textbf{10.} $v_1 = \hbar/m_ea_0 = \qty{2.19e6}{m/s}$~; $v_1/c = 7.3 \times 10^{-3} =
1/137$.

\textbf{11.} $2 \to 1$~: \qty{10.2}{eV}, \qty{122}{nm} (ultraviolet)~; $3 \to 2$~:
\qty{1.89}{eV}, \qty{656}{nm} (rouge)~; $4 \to 2$~: \qty{2.55}{eV}, \qty{486}{nm}
(bleu-vert)~: les raies de Balmer sont les raies visibles.

\textbf{12.} $hc/\lambda = |E_1|(1/n^2 - 1/p^2)$~: $R_H = |E_1|/hc = 13.6/1240 =
\qty{1.097e-2}{nm^{-1}} = \qty{1.097e7}{m^{-1}}$. Limite de Balmer ($p \to \infty$)~:
$\lambda = 4/R_H = \qty{365}{nm}$.

\textbf{13.} \qty{13.6}{eV}~; $\lambda = 1240/13.6 = \qty{91}{nm}$.

\textbf{14.} $v_n = v_1/n$, $\lambda_n = h/m_ev_n = nh/m_ev_1 = 2\pi na_0$ (puisque $m_ev_1a_0
= \hbar$)~; $2\pi r_n = 2\pi n^2a_0 = n\lambda_n$~: l'orbite porte $n$ ondes
entières --- la règle de Bohr est une règle d'\omterm{thm:b1:wave-propagation:standing}{onde stationnaire}.

\textbf{15.} Faux~: l'électron n'a pas d'orbite (il rayonnerait~; l'\omterm{prop:b1:quantum-introduction:box}{état
fondamental} a un \omterm{def:b1:angular-momentum:L}{moment cinétique} nul), et le modèle échoue dès qu'un atome
a deux électrons. Juste~: les niveaux de l'hydrogène, donc chaque raie de son
spectre. L'orbite est remplacée par une \omterm{def:b1:quantum-introduction:wavefunction}{fonction d'onde}, un nuage de
probabilité de taille $\sim n^2a_0$.

\textbf{16.} \omterm{thm:b1:wave-propagation:standing}{Ondes stationnaires}~: $L = n\lambda/2$, $p_n = h/\lambda_n = nh/2L$, $E_n =
p_n^2/2m = n^2h^2/8mL^2$. La formule de la sphère est la même avec $L \to R$~:
l'\omterm{prop:b1:quantum-introduction:box}{état fondamental} loge une demi-longueur d'onde dans le rayon.

\textbf{17.} $E_e = 6.02 \times 10^{-38}/(0.13 \times 4 \times 10^{-18}) = \qty{0.72}{eV}$~;
$E_h = 0.13/0.45 \times 0.72 = \qty{0.21}{eV}$.

\textbf{18.} $R = \qty{2}{nm}$~: $1.74 + 0.93 = \qty{2.67}{eV}$, \qty{464}{nm}, bleu~;
\qty{3}{nm}~: \qty{2.15}{eV}, \qty{577}{nm}, jaune~; \qty{4}{nm}~: \qty{1.97}{eV},
\qty{629}{nm}, rouge.

\textbf{19.} L'énergie de confinement croît comme $1/R^2$~: une boîte plus petite
relève les niveaux, le photon est plus énergétique, la lumière plus bleue.

\textbf{20.} $\Delta p \sim \hbar/R = \qty{5.3e-26}{kg.m/s}$, $E \sim \Delta p^2/2m_e^\ast =
\qty{1.2e-20}{J} = \qty{0.07}{eV}$~: le bon ordre (le coefficient exact
$\pi^2/2$ et l'estimation grossière de $\Delta p$ expliquent le facteur dix).

\textbf{21.} $1240/577 = \qty{2.15}{eV} = \qty{3.4e-19}{J}$~: $10^{-9}/3.4 \times 10^{-19}
= 2.9 \times 10^9$ photons par seconde.

\textbf{22.} $10^{-3}/(1.96 \times 1.6 \times 10^{-19}) = 3.2 \times 10^{15}$ par seconde.

\textbf{23.} $100 \times 2.48 \times 1.6 \times 10^{-19} = \qty{4e-17}{J}$ en \qty{0.1}{s}~:
\qty{4e-16}{W}.

\textbf{24.} $h\nu = \qty{6.6e-26}{J}$~: $1.5 \times 10^{13}$ photons par seconde ---
et chacun est $10^5$ fois plus petit que l'énergie thermique $k_BT$ de
l'antenne~: aucun espoir, et aucun besoin~; la description ondulatoire y est exacte.

\textbf{25.} La taille des atomes, $a_0 = \qty{53}{pm}$~; leur échelle d'énergie,
\qty{13.6}{eV} (d'où des photons de quelques eV, la chimie, le spectre
visible)~; et la vitesse de leurs électrons, $\alpha c \approx \qty{2e6}{m/s}$ ---
les trois à partir de $h$, $e$ et $m_e$ (et $\varepsilon_0$), sans rien ajuster.
\end{solution}
